\section{Predication as the seed out of potential}

\begin{proposition}[Entity and predicate]\label{thm:entity-predicate}\lean{Infomaterial.Entity}\leanok\uses{def:mark, thm:kind}
An \emph{entity} is a mark in identity-carrying role, and a
\emph{predicate} is a mark in content-carrying role. At this base both are
abbreviations for \texttt{Mark}; typed predicate structure enters later
through positing and valuation.
\end{proposition}

\begin{proposition}[The predication cell]\label{thm:predication}\lean{Infomaterial.PredicationCell}\leanok\uses{thm:entity-predicate, thm:articulation-cell-corner}
The predication cell is the oriented triple $(s,p,o)$: subject entity,
predicate mark, and object entity. Orientation belongs to the datum, so the
same predicate read from opposite endpoints is a different cell off the
diagonal.
\end{proposition}

\begin{proposition}[The reflexive seed]\label{thm:predication-seed}\lean{Infomaterial.PredicationCell.seed}\leanok\uses{thm:predication, def:utterance}
The seed of an utterance with token $t$ and claim $c$ is the reflexive
predication cell $(t,c,t)$. Predicating first creates the diagonal assertion
position before later constructions move off the diagonal.
\end{proposition}

\begin{proposition}[Predicating proper]\label{thm:predicating-proper}\lean{Infomaterial.PredicationCell.predicating}\leanok\uses{thm:predication-seed, thm:articulation-cell-initial}
Predicating is defined on an articulation cell only with a proof that the
cell is in the potential corner. Under that incidence condition it returns
the seed of the cell's utterance.
\end{proposition}
